% ---------------------------------------------------------------------------
% Chromatic numbers of the plane with forbidden distances:
% a seven-color lower bound for {1,(sqrt6+sqrt2)/2} and smaller
% six-chromatic two-distance graphs.
% ---------------------------------------------------------------------------
\documentclass[11pt]{article}

\usepackage[margin=1.1in]{geometry}
\usepackage{amsmath,amssymb,amsthm}
\usepackage{booktabs}
\usepackage{array}
\usepackage{enumitem}
\usepackage{xcolor}
\usepackage{xurl}
\usepackage[colorlinks=true,linkcolor=blue!50!black,citecolor=blue!50!black,urlcolor=blue!50!black]{hyperref}

\newtheorem{theorem}{Theorem}[section]
\newtheorem{lemma}[theorem]{Lemma}
\newtheorem{proposition}[theorem]{Proposition}
\newtheorem{corollary}[theorem]{Corollary}
\newtheorem{conjecture}[theorem]{Conjecture}
\theoremstyle{definition}
\newtheorem{definition}[theorem]{Definition}
\newtheorem{question}[theorem]{Question}
\newtheorem{problem}[theorem]{Problem}
\theoremstyle{remark}
\newtheorem{remark}[theorem]{Remark}

\newcommand{\R}{\mathbb{R}}
\newcommand{\Z}{\mathbb{Z}}
\newcommand{\Q}{\mathbb{Q}}
\newcommand{\E}{\mathbb{E}}
\newcommand{\dd}{d}                          % the second distance
\newcommand{\chid}[1]{\chi\bigl(\R^2;#1\bigr)}
\DeclareMathOperator{\Iso}{Isom}

\title{Chromatic Numbers of the Plane with Two Forbidden Distances:\\
a Seven-Color Lower Bound and Smaller Six-Chromatic Graphs}

\author{Henry Bass}

\date{October 2026}

\begin{document}
\maketitle

\begin{abstract}
Let $\chid{D}$ denote the minimum number of colors required to color the
plane so that no two points whose distance lies in $D$ receive the same
color, and set $\dd=(\sqrt6+\sqrt2)/2=2\cos(\pi/12)$, so that
$\dd^2=2+\sqrt3$ and $\dd^2-1=\sqrt2\,\dd=1+\sqrt3$.  Our main result is
\[
  \chid{\{1,\dd\}}\;\ge\;7
\]
for arbitrary colorings, with no measurability assumed (the averaging uses an invariant mean on the amenable group $E(2)$, a theorem of ZFC): for every proper
$6$-coloring, invariant monochromatic-pair densities $p_{\sqrt2}$ and
$p_{1+\sqrt3}$ would have to satisfy $p_{\sqrt2}+p_{1+\sqrt3}\ge 1$
(an $8$-point pigeonhole gadget of Parts) and $p_{\sqrt2},p_{1+\sqrt3}\le
5/11$ (two explicit $12$-point chain configurations whose long-range pairs
land exactly on the forbidden distances via the identity
$\dd^2-1=\sqrt2\,\dd$), a contradiction since $1\le 10/11$ is false.  The
argument realizes a 2019 sketch of J.~Parts in a Polymath16 blog comment;
the contribution here is a fully explicit finite form with an explicit
gap ($p_r\le 5/11=\tfrac12-\tfrac1{22}$), each of the three finite
statements carrying a DRAT proof trace checked by \texttt{drat-trim},
and the full $32$-point geometry certified in exact symbolic arithmetic
over explicit quadratic towers.  The argument does not use the
measurable-to-arbitrary transfer of OpenAI's 2026 proof that
$\chi(\R^2)\ge6$.

We also give, with exact-arithmetic edge sets and DRAT-certified
non-colorability: a $26$-vertex $6$-chromatic $\{1,\dd\}$-graph (to our
knowledge the smallest $6$-chromatic two-distance graph known for any
ratio, improving $31$ [Parts 2020] and, for this ratio, $117$ [De~Neve
et al.\ 2025]); a $30$-vertex $6$-chromatic $\{1,\sqrt3\}$-graph
(previously $33$); the six-coloring continuum
$X_6\supseteq[0.354,0.667504532081]$ (previously $[0.354,0.657]$), by
$59$ exact rational tilings; and the edge bound $e_5\le 2278$ for
$5$-chromatic unit-distance graphs, via a $2278$-edge non-$4$-colorable
spanning subgraph of the $509$-vertex record graph (previously $2406$).
As by-products we record forced monochromatic densities
$p_{\sqrt3}\ge 1/28$, $p_{\dd}\ge 1/37$ and $p_{1/\sqrt3}\ge1/67$ for
invariant unit-distance $5$-colorings, together with elementary
color-count-free density caps such as $p_{1/\sqrt3}\le1/3$; since
$\chi(\R^2)\ge6$ is now known, these statements about $5$-colorings are
vacuous, and we keep them as finite combinatorial statements.  All
certificates and verifiers are included in the accompanying
\texttt{certificates/} folder (Appendix~\ref{app:artifacts}).
\end{abstract}

\section{Introduction}\label{sec:intro}

The Hadwiger--Nelson problem asks for the chromatic number of the plane,
$\chi(\R^2)=\chid{\{1\}}$: the least number of colors so that no two points
at distance $1$ share a color.  After more than half a century at
$4\le\chi(\R^2)\le 7$, de~Grey \cite{deGrey18} proved $\chi(\R^2)\ge 5$
with a $1581$-vertex unit-distance graph, launching the Polymath16 project
\cite{Polymath16} and a rapid sequence of smaller witnesses
\cite{Heule18,ExooIsmailescu18,Parts20b}.

In the \emph{two-distance} variant one forbids two distances
simultaneously: $\chid{\{1,\dd\}}$ is the chromatic number of the graph on
$\R^2$ with edges between points at distance $1$ \emph{or} $\dd$.  Finite
$6$-chromatic witnesses were known for exactly four ratios: the golden
ratio $(1+\sqrt5)/2$ (Huddleston \cite{OTH08}, without an explicit
graph; an explicit $126$-vertex graph was given by de~Grey in
Polymath16 \cite{Polymath16}, and the $31$-vertex record is due to Parts
\cite{Parts20a}), $\dd=2$
(Exoo--Ismailescu \cite{ExooIsmailescu20}, reduced to $397$ vertices by
Parts \cite{Parts20a}), and $\dd=\sqrt3$ and $\dd=(\sqrt6+\sqrt2)/2$
(De~Neve, Vanden Kerchove, Colle, Tavernier, and Pickavet
\cite{DeNeve25}, with $33$ and $117$ vertices respectively; the lower
bound $6$ for these two ratios had been obtained earlier by
P\'alv\"olgyi and \'Agoston with the probabilistic method
\cite{PalvolgyiAgoston19}).  For no finite set $D$ of distances was a
bound $\chid{D}\ge 7$ published, although J.~Parts sketched a
probabilistic argument for exactly this in a 2019 Polymath16 comment
\cite{Parts19}.

\emph{Note added (October 2026).} The results of this paper were
obtained in August 2026.  Since then OpenAI~\cite{OAI} proved
$\chi(\R^2)\ge6$, so every two-distance chromatic number of the plane is
at least $6$, the statements about $5$-colorings in
Sections~\ref{sec:forced} and~\ref{sec:k5caps} are vacuous, and the win
condition of Section~\ref{sec:k5caps} is moot.  The explicit small
$6$-chromatic graphs of Section~\ref{sec:record} are unaffected.  The
transfer principle of~\cite{OAI} extends to countable sets of algebraic
distances (Crocker~\cite{Crocker}, conditionally on a theorem
of~\cite{OAI}, and the companion manuscript~\cite{BassR3}, with a Lean
formalization), giving $\chid{D}\ge 1/m_D(\R^2)$, where $m_D$ is the
maximal density of a measurable set avoiding $D$.  The companion
manuscript~\cite{BassMono} combines this with Fourier density
certificates and obtains, among other things, $\chid{\{1,\dd'\}}\ge7$
for every algebraic $\dd'\in[1.77,1.92]\cup[7,\infty)$; with Bukh's
theorem~\cite{Bukh} one also gets $\chid{\{1,N\}}\ge17$ for all
sufficiently large algebraic $N$, with no explicit threshold.  The value
$\dd=(\sqrt6+\sqrt2)/2\approx1.932$ lies outside the ranges covered
there, and Theorem~\ref{thm:main} is proved differently: by three finite
SAT-certified configurations and an invariant mean, without the
transfer.  Euclidean Ramsey consequences of the transfer are developed
in~\cite{BassRamsey}.

All searches and certificates described below were computer-assisted.
Throughout, $\dd:=(\sqrt6+\sqrt2)/2=2\cos(\pi/12)\approx
1.93185$, whose relevant algebra is
\begin{equation}\label{eq:identity}
  \dd^{\,2}=2+\sqrt3,
  \qquad
  \dd^{\,2}-1 \;=\; 1+\sqrt3 \;=\; \sqrt2\,\dd .
\end{equation}

\begin{theorem}[Main theorem; invariant-coloring framework]\label{thm:main}
$\chid{\{1,\dd\}}\ge 7$.  More precisely, there is no proper $6$-coloring
$c\colon\R^2\to\{1,\dots,6\}$ avoiding monochromatic pairs at distances
$1$ and $\dd$: for any such $c$, invariant monochromatic-pair densities
$p_{\sqrt2}(c)$ and $p_{1+\sqrt3}(c)$ exist (via an invariant mean on the
isometry group) and satisfy the incompatible system
\[
  p_{\sqrt2}+p_{1+\sqrt3}\ \ge\ 1,
  \qquad
  p_{\sqrt2}\ \le\ \tfrac{5}{11},
  \qquad
  p_{1+\sqrt3}\ \le\ \tfrac{5}{11}.
\]
\end{theorem}

The three displayed inequalities come from three finite configurations
($8+12+12=32$ points in total), each certified by a DRAT proof trace
(\texttt{certificates/seven-colors/}; Section~\ref{sec:seven} and
Appendix~\ref{app:artifacts}).  The framework, its provenance in ZFC, and
its limitations are discussed in
Sections~\ref{sec:prelim}--\ref{sec:seven}; a counting blueprint that
would replace the framework by a single finite graph is given in
Section~\ref{sec:blueprint}.

The remaining results rest on finite statements with exact-arithmetic
and DRAT certificates:

\begin{theorem}\label{thm:records}
\leavevmode
\begin{enumerate}[label=(\alph*)]
\item There are $26$ points in the plane spanning $85$ pairs at distance
exactly $1$ or exactly $\dd$ whose $\{1,\dd\}$-graph has chromatic number
$6$ and is vertex-critical.  To our knowledge this is the smallest
$6$-chromatic two-distance graph known for any distance pair (previously $31$
\cite{Parts20a}; for this pair, $117$ \cite{DeNeve25}).
\item There are $30$ points spanning $100$ pairs at distance exactly $1$
or exactly $\sqrt3$ whose graph has chromatic number $6$ and is
vertex-critical (previously $33$ \cite{DeNeve25}).
\item The record $509$-vertex $5$-chromatic unit-distance graph
\cite{Parts20b,deGreyParts22} contains a spanning subgraph with $2278$
edges that is not $4$-colorable.  Hence $e_5\le 2278$ (previously
$e_5\le 2406$ \cite{deGreyParts22}).
\item For every proper unit-distance $5$-coloring of the plane and every
invariant mean, the monochromatic-pair densities of
Definition~\ref{def:pr} satisfy $p_{\sqrt3}\ge 1/28$, $p_{\dd}\ge 1/37$
and $p_{1/\sqrt3}\ge1/67$.  (Since $\chi(\R^2)\ge6$~\cite{OAI}, no such
coloring exists; the content is the underlying finite statement, see
Section~\ref{sec:forced}.)
\end{enumerate}
\end{theorem}

Each item is backed by files in \texttt{certificates/}: (a) by
\texttt{g26/}, (b) by \texttt{g30/}, (c) by \texttt{e5-2278/}, and (d)
by \texttt{g30/}, \texttt{g26/} and \texttt{forced-density/}
respectively.

\begin{center}
\small
\begin{tabular}{@{}p{4.9cm}p{2.8cm}lp{3.5cm}@{}}
\toprule
Quantity & Previous & Here & Certificate \\
\midrule
Smallest $6$-chromatic two-distance graph, any ratio &
$31$ \cite{Parts20a} & $\mathbf{26}$ & exact edges $+$ DRAT
(\texttt{g26/}) \\
\dots\ for $\dd=(\sqrt6+\sqrt2)/2$ &
$117$ \cite{DeNeve25} & $\mathbf{26}$ & exact edges $+$ DRAT
(\texttt{g26/}) \\
\dots\ for $\dd=\sqrt3$ &
$33$ \cite{DeNeve25} & $\mathbf{30}$ & exact edges $+$ DRAT
(\texttt{g30/}) \\
$\chid{\{1,(\sqrt6+\sqrt2)/2\}}$ &
$\ge 6$ \cite{PalvolgyiAgoston19,DeNeve25}; $\ge 7$ sketched
\cite{Parts19} &
$\mathbf{\ge 7}$ & exact geometry $+$ 3 DRAT
(\texttt{seven-colors/}) \\
Edges of the $509$-vertex $5$-chromatic UD graph &
$\le 2406$ \cite{deGreyParts22} & $\mathbf{\le 2278}$ & DRAT
(\texttt{e5-2278/}) \\
Six-coloring continuum $X_6$ &
$[0.354,0.657]$ \cite{MPSZ} & $[0.354,\mathbf{0.6675045}\dots]$ &
59 exact rational tilings (\texttt{six-coloring-}\linebreak[0]\texttt{continuum/}) \\
Forced densities $p_{\sqrt3}$, $p_{\dd}$, $p_{1/\sqrt3}$ (vacuous
since~\cite{OAI}) &
--- & $\mathbf{1/28}$, $\mathbf{1/37}$, $\mathbf{1/67}$ & DRAT
(\texttt{g30/}, \texttt{g26/}, \texttt{forced-density/}) \\
\bottomrule
\end{tabular}
\end{center}

\paragraph{Attribution.}
The $\chi\ge7$ argument is, in its skeleton, due to Jaan Parts, who
sketched it on 24 December 2019 in comment~24686 on the fifteenth
Polymath16 thread \cite{Parts19} and, to our knowledge, never published
it.  Our contribution is (i) fully explicit finite certificates --- the
specific $8$- and $12$-point configurations, their exact geometry, and
their SAT/DRAT verification --- (ii) the explicit constant $5/11$ for
the upper bound on $p_r$, and (iii) a counting blueprint toward a
choice-free finite witness (Section~\ref{sec:blueprint}).  The record
graphs of Theorem~\ref{thm:records}(a,b) were seeded from the published
coordinates of De~Neve et al.\ \cite{DeNeve25}.

\paragraph{Organization.}
Section~\ref{sec:prelim} sets up two-distance graphs and the
invariant-coloring framework.  Section~\ref{sec:seven} proves
Theorem~\ref{thm:main}.  Section~\ref{sec:blueprint} describes the
counting blueprint toward a fully finite witness, together with
computational observations about it.  Section~\ref{sec:record} presents
the record graphs with complete coordinate tables and verification
recipes.  Section~\ref{sec:ud} contains the unit-distance results
(Theorem~\ref{thm:records}(c,d)).  Section~\ref{sec:k5caps} develops
finite caps valid for any number of colors, a bipartiteness obstruction
at the distances $\sqrt3$ and $2$, and the (now moot) win condition for
$\chi(\R^2)\ge6$.  Section~\ref{sec:continuum} extends the six-coloring
continuum.  Section~\ref{sec:open} collects open problems.
Appendix~\ref{app:artifacts} documents the artifact bundle.

\section{Preliminaries}\label{sec:prelim}

\subsection{Two-distance graphs}

For a finite $P\subset\R^2$ and a set $D$ of positive reals, let $G(P;D)$
be the graph on $P$ with $uv$ an edge iff $|u-v|\in D$.  We write
$\chid{D}$ for the chromatic number of the (infinite) graph $G(\R^2;D)$
and abbreviate $G(P)=G(P;\{1,\dd\})$ when $D=\{1,\dd\}$ is understood.
By the de~Bruijn--Erd\H{o}s theorem (with choice), $\chid{D}\ge k$
holds iff some finite $P$ has $\chi(G(P;D))\ge k$; part of the interest
of Theorem~\ref{thm:main} is precisely that it does \emph{not} yet come
with such a finite witness (see Section~\ref{sec:caveat}).

All configurations in Section~\ref{sec:record} have coordinates in a real
quadratic-tower field ($\Q(\sqrt3,\sqrt5)$ or $\Q(\sqrt3,\sqrt7)$) with
denominators dividing $8$, so squared distances are computed in exact
integer arithmetic and edge membership is decidable with no floating
point.  The $\chi\ge7$ configurations of Section~\ref{sec:seven} are
likewise exact algebraic configurations, certified end-to-end by
symbolic computation in explicit quadratic towers
(Section~\ref{sec:exact}).

\subsection{The invariant-coloring framework}\label{sec:framework}

We use the probabilistic formulation of Hadwiger--Nelson-type problems
developed by the Polymath16 project \cite{Polymath16}, \'Agoston
\cite{Agoston21}, and Gwyn--Stavrianos \cite{GwynStavrianos20}.  The
isometry group $E(2)=\Iso(\R^2)$ is solvable, hence amenable, so it
admits a left-invariant mean $\Lambda$ on $\ell^\infty(E(2))$ (this uses
Hahn--Banach; it is a theorem of ZFC).

\begin{definition}\label{def:pr}
Fix any proper coloring $c$ of $G(\R^2;D)$ with $k$ colors (no
measurability assumed) and fix an invariant mean $\Lambda$.  For $r>0$
define the \emph{monochromatic-pair density}
\[
  p_r \;=\; p_r(c,\Lambda)\;=\;
  \Lambda_{g\in E(2)}\Bigl[\,\mathbf 1\bigl(c(g\cdot(0,0))=c(g\cdot(r,0))\bigr)\Bigr]
  \ \in[0,1].
\]
\end{definition}

By left-invariance, for \emph{every} ordered pair $(x,y)$ with $|x-y|=r$
we have $\Lambda_g[\mathbf 1(c(gx)=c(gy))]=p_r$.  This yields the
transfer principle we use throughout:

\begin{lemma}[Finite transfer]\label{lem:transfer}
Let $Q\subset\R^2$ be finite with designated pairs
$(a_1,b_1),\dots,(a_m,b_m)$, $|a_i-b_i|=r_i$, and let $t_{\min}$
(resp.\ $t_{\max}$) be the minimum (resp.\ maximum), over all proper
$k$-colorings of the finite graph $G(Q;D)$, of the number of
monochromatic designated pairs.  Then for every proper $k$-coloring $c$
of $G(\R^2;D)$ and every invariant mean $\Lambda$,
\[
  t_{\min}\ \le\ \sum_{i=1}^{m} p_{r_i}(c,\Lambda)\ \le\ t_{\max}.
\]
\end{lemma}

\begin{proof}
For each $g\in E(2)$, the restriction of $c$ to $gQ$ is a proper
$k$-coloring of $G(gQ;D)\cong G(Q;D)$, so the number of monochromatic
designated pairs of $gQ$ lies in $[t_{\min},t_{\max}]$ pointwise in $g$.
Apply $\Lambda$ and use linearity together with the paragraph above,
which gives $\Lambda_g[\mathbf 1(c(ga_i)=c(gb_i))]=p_{r_i}$.
\end{proof}

Only one-sided bounds are needed in applications: a lower bound
$t\le t_{\min}$ or an upper bound $t\ge t_{\max}$.  Each bound we use is
certified by refuting the decision statement ``a proper $k$-coloring
with $\le t-1$ (resp.\ $\ge t+1$) monochromatic designated pairs
exists'': the statement is encoded as a CNF formula, and its
unsatisfiability is certified by a DRAT proof trace checked by
\texttt{drat-trim} \cite{drattrim}.  (The optima were located with an
exact MaxSAT solver, RC2 over CaDiCaL, and are attained by explicit
colorings, but only the DRAT-certified direction enters the proofs.)

\subsection{Prior work used directly}

De~Grey's graph \cite{deGrey18} and its successors
\cite{Heule18,ExooIsmailescu18,Parts20b} give $\chi(\R^2)\ge5$; the
$509$-vertex record and its $2406$-edge trim are from
\cite{Parts20b,deGreyParts22}.  The two-distance $6$-chromatic graphs we
improve are from \cite{Parts20a,ExooIsmailescu20,DeNeve25}.  The
probabilistic framework is from \cite{Polymath16,Agoston21}, and the
finite-graph transfer of badness bounds from \cite{GwynStavrianos20}.
The $\chi\ge7$ sketch is Parts' \cite{Parts19}.  Set-theoretic
sensitivity of chromatic numbers of distance graphs is due to
Shelah--Soifer \cite{ShelahSoifer03}; general background is in Soifer
\cite{Soifer09}.

\section{Seven colors for two distances}\label{sec:seven}

Throughout this section $k=6$, $D=\{1,\dd\}$, and $c$ is an arbitrary
proper $6$-coloring of $G(\R^2;\{1,\dd\})$, assumed to exist for
contradiction; $p_r=p_r(c,\Lambda)$ for a fixed invariant mean.  The
two auxiliary distances are $\sqrt2$ and $1+\sqrt3$ --- not themselves
forbidden, but tied to $\{1,\dd\}$ by the identity \eqref{eq:identity}.

\subsection{The V8 pigeonhole}\label{sec:v8}

Parts' \emph{V8 gadget} consists of two concentric squares: the unit
square $I=\{(\pm\tfrac12,\pm\tfrac12)\}$ and the outer square
$O=\{(\pm s,0),(0,\pm s)\}$ with $s=\dd/\sqrt2=(1+\sqrt3)/2$.

\begin{lemma}\label{lem:v8geom}
$G(I\cup O;\{1,\dd\})=K_8$ minus the perfect matching consisting of the
four ``diagonals'': the two inner diagonals at distance $\sqrt2$ and the
two outer diagonals at distance $2s=1+\sqrt3$.
\end{lemma}

\begin{proof}
Direct computation.  Inner sides have length $1$; inner diagonals
$\sqrt2$.  Outer sides have squared length $2s^2=2+\sqrt3=\dd^2$; outer
diagonals $2s=1+\sqrt3$.  For an inner corner and an outer vertex, e.g.\
$(\tfrac12,\tfrac12)$ and $(s,0)$: the difference is
$(\tfrac{\sqrt3}2,\tfrac12)$ of length $1$; and against $(-s,0)$ the
squared distance is $\bigl(\tfrac{2+\sqrt3}{2}\bigr)^2+\tfrac14
=2+\sqrt3=\dd^2$.  By symmetry each inner vertex is at distance $1$ from
the two nearest outer vertices and $\dd$ from the two farthest.  Thus all
$24$ pairs except the four diagonals are edges.
\end{proof}

\begin{proposition}[V8 pigeonhole]\label{prop:v8}
Every proper $6$-coloring of $G(I\cup O;\{1,\dd\})$ makes at least $2$ of
the four diagonal pairs monochromatic.  The bound is exact, and all $11$
diagonal patterns of weight $\ge2$ are realizable.  Certificate:
\texttt{certificates/seven-colors/v8cert.cnf} (proper $6$-coloring
$\wedge$ at most one monochromatic diagonal; $84$ variables, $267$
clauses) is UNSAT, with DRAT trace \texttt{v8cert.drat} verified by
\texttt{drat-trim}.
\end{proposition}

\begin{proof}
By Lemma~\ref{lem:v8geom} every color class is either a single vertex or
one of the four diagonal pairs.  If exactly $j$ diagonals are
monochromatic, the coloring uses $8-j$ colors, so $8-j\le6$ forces
$j\ge2$; conversely, for any $j\ge2$ diagonals, coloring each of them
with its own color and the remaining $8-2j$ vertices with distinct
further colors uses $8-j\le6$ colors.  The DRAT certificate checks the
first statement mechanically.
\end{proof}

By Lemma~\ref{lem:transfer} with the four diagonals as
designated pairs (two at $\sqrt2$, two at $1+\sqrt3$) and $t_{\min}=2$:

\begin{corollary}\label{cor:v8}
$2p_{\sqrt2}+2p_{1+\sqrt3}\ \ge\ 2$, i.e.\
$p_{\sqrt2}+p_{1+\sqrt3}\ \ge\ 1$.
\end{corollary}

\subsection{The chain configurations and the exact identity}\label{sec:chains}

The counterweight to Corollary~\ref{cor:v8} is a pair of $12$-point
configurations bounding $p_{\sqrt2}$ and $p_{1+\sqrt3}$ from above.
Fix a radius $r\in\{\sqrt2,\,1+\sqrt3\}$ and let
\[
  \varphi_1=2\arcsin\frac{1}{2r},\qquad
  \varphi_{\dd}=2\arcsin\frac{\dd}{2r}
\]
be the angles subtending chords of length $1$ and $\dd$ on a circle of
radius $r$; equivalently $\cos\varphi_1=1-\tfrac1{2r^2}$ and
$\cos\varphi_{\dd}=1-\tfrac{\dd^2}{2r^2}$.  Concretely,
\[
  r=\sqrt2:\ \ \cos\varphi_1=\tfrac34,\ \
  \cos\varphi_{\dd}=\tfrac{2-\sqrt3}{4};
  \qquad
  r=1+\sqrt3:\ \ \cos\varphi_1=\tfrac{2+\sqrt3}{4},\ \
  \cos\varphi_{\dd}=\tfrac34,
\]
a pleasing duality: the two radii exchange the roles of the two cosines.

Set $A=\varphi_{\dd}-\varphi_1$ and $B=\varphi_{\dd}+\varphi_1$, and
define the \emph{mixed chain} $X_0,\dots,X_6$ by $X_0=(0,0)$ and
$X_{i+1}=X_i+r(\cos h_i,\sin h_i)$, where the heading sequence has period
four:
\[
  (h_0,h_1,h_2,h_3)\;=\;\bigl(0,\ \pi-A,\ -(A+B),\ \pi-B\bigr),
  \qquad h_{i+4}=h_i .
\]
Thus consecutive chain points are at distance exactly $r$.  The key fact
is that the four-step displacement lands \emph{exactly} on a forbidden
distance:

\begin{lemma}[Four-step closure]\label{lem:closure}
For all $i$, $X_{i+4}-X_i$ is the same vector, of length
$(\dd^2-1)/r$.  Consequently, by \eqref{eq:identity}:
for $r=\sqrt2$, $|X_{i+4}-X_i|=(1+\sqrt3)/\sqrt2=\dd$;
for $r=1+\sqrt3$, $|X_{i+4}-X_i|=1$.
\end{lemma}

\begin{proof}
Identifying $\R^2$ with $\mathbb C$,
$X_{i+4}-X_i=r\sum_{j=i}^{i+3}e^{\mathrm i h_j}
=r\sum_{j=0}^{3}e^{\mathrm i h_j}$, since the headings have period four.
Now
\[
  \sum_{j=0}^{3}e^{\mathrm i h_j}
  \;=\;1-e^{-\mathrm iA}+e^{-\mathrm i(A+B)}-e^{-\mathrm iB}
  \;=\;\bigl(1-e^{-\mathrm iA}\bigr)\bigl(1-e^{-\mathrm iB}\bigr),
\]
whose modulus is $4\sin\tfrac A2\sin\tfrac B2
=2\bigl(\cos\tfrac{B-A}2-\cos\tfrac{B+A}2\bigr)
=2(\cos\varphi_1-\cos\varphi_{\dd})
=2\cdot\tfrac{\dd^2-1}{2r^2}=\tfrac{\dd^2-1}{r^2}$.
Multiplying by $r$ gives $|X_{i+4}-X_i|=(\dd^2-1)/r$.
\end{proof}

So for $r=\sqrt2$ the pairs $(X_i,X_{i+4})$, $i=0,1,2$, are genuine
$\dd$-edges of the two-distance graph, and for $r=1+\sqrt3$ they are unit
edges: the chain's long-range structure is policed by the forbidden
distances even though its steps are not.  To stiffen the configuration we
add five \emph{apexes} $C_1,\dots,C_5$: $C_i$ is the rotation of
$X_{i-1}$ about $X_i$ through $\pm\varphi_1$ (one branch per $i$).  Then
$|C_iX_i|=r$ and $|C_iX_{i-1}|=1$ automatically, and the branch is chosen
so that $|C_iX_{i+1}|=\dd$ --- exactly, for every $i$, at both radii
(Section~\ref{sec:exact}); the apexes are distinct from all chain
points.  The resulting configuration
$Q_r=\{X_0,\dots,X_6,C_1,\dots,C_5\}$ has $12$ points; its
$\{1,\dd\}$-graph has $14$ edges ($5$ unit and $9$ at distance $\dd$ for
$r=\sqrt2$; $9$ and $5$ respectively for $r=1+\sqrt3$).  The
\emph{charged slots} are the eleven designated pairs at distance $r$:
\[
  (X_i,X_{i+1})\ \ (0\le i\le 5)
  \qquad\text{and}\qquad
  (X_i,C_i)\ \ (1\le i\le 5).
\]

\begin{proposition}[Chain caps]\label{prop:caps}
For $r=\sqrt2$ (chain $A$) and for $r=1+\sqrt3$ (chain $B$): in every
proper $6$-coloring of $G(Q_r;\{1,\dd\})$, at most $5$ of the $11$
charged slots are monochromatic, and $5$ is attained.  Certificates (in
\texttt{certificates/seven-colors/}): \texttt{capA.cnf} and
\texttt{capB.cnf} (proper $6$-coloring $\wedge$ at least $6$
monochromatic slots; $188$ variables, $646$ clauses each) are UNSAT with
DRAT traces \texttt{capA.drat}, \texttt{capB.drat} verified by
\texttt{drat-trim}; independently, \texttt{verify\_exact.py} re-proves
the cap by exhaustive enumeration of slot subsets on the \emph{exact}
edge lists of Section~\ref{sec:exact} and exhibits a coloring attaining
$5$.
\end{proposition}

By Lemma~\ref{lem:transfer} with the eleven slots (all at distance $r$)
and $t_{\max}=5$:

\begin{corollary}\label{cor:caps}
$11\,p_{\sqrt2}\le 5$ and $11\,p_{1+\sqrt3}\le 5$; that is,
$p_{\sqrt2}\le \tfrac5{11}$ and $p_{1+\sqrt3}\le\tfrac5{11}$.
\end{corollary}

Thus a single $12$-point configuration forces slot mono-density
$\le\tfrac5{11}=\tfrac12-\tfrac1{22}$, and $5/11$ is optimal for this
configuration.  Parts' sketch \cite{Parts19} instead obtains a bound
$p_r<\tfrac12$ by counting along long zigzag chains, without an explicit
constant.

\subsection{Exact-arithmetic certification of the chains}\label{sec:exact}

Every geometric statement about $Q_{\sqrt2}$ (chain $A$) and
$Q_{1+\sqrt3}$ (chain $B$) --- including the apex placements --- has been
verified in exact symbolic arithmetic; no floating point enters the
$\chi\ge7$ certificates.  The verifier is
\texttt{certificates/seven-colors/verify\_exact.py} ($\sim$250 lines, on
\texttt{sympy} $+$ \texttt{python-sat}); its design:

\begin{itemize}
\item \emph{Towers.}  All coordinates of chain $A$ lie in
$K_A=\Q(\sqrt2,\sqrt3,\sqrt7,S)$ with $S=\sqrt{9+4\sqrt3}$
(note $\sin\varphi_{\dd}=S/4$, $\sin\varphi_1=\sqrt7/4$ at $r=\sqrt2$),
a tower of nominal degree $16$; those of chain $B$ lie in
$K_B=\Q(\sqrt3,\sqrt7,T)$ with $T=\sqrt{9-4\sqrt3}$
(here $\sin\varphi_1=T/4$, $\sin\varphi_{\dd}=\sqrt7/4$), of degree $8$.
Elements are polynomials in the generators, reduced to normal form
modulo the defining relations $S^2=9+4\sqrt3$, $w_2^2=2$, $w_7^2=7$,
$w_3^2=3$ (resp.\ $T^2=9-4\sqrt3$, \dots).  Under a lexicographic order
the relation set has pairwise-coprime pure-square leading monomials,
hence is a Gr\"obner basis, so normal forms are unique and the quotient
ring has dimension $16$ (resp.\ $8$).
\item \emph{Soundness certificate.}  Nonzero normal form certifies a
nonzero field element only if the quotient ring \emph{is} the field.
This is certified separately: the minimal polynomial of the primitive
element $\sqrt2+\sqrt3+\sqrt7+\sqrt{9+4\sqrt3}$ (resp.\
$\sqrt3+\sqrt7+\sqrt{9-4\sqrt3}$) is computed and its degree confirmed
to be $16$ (resp.\ $8$).  Every subsequent equality/inequality decision
is therefore exact and unconditional.
\item \emph{Checks.}  For each chain: the headings are exact unit
vectors; all six steps satisfy $|X_iX_{i+1}|^2=r^2$ exactly; the five
apexes satisfy $|C_iX_i|^2=r^2$, $|C_iX_{i-1}|^2=1$, and
$|C_iX_{i+1}|^2=\dd^2=2+\sqrt3$ exactly; the three long pairs satisfy
$|X_iX_{i+4}|^2=(\dd^2-1)^2/r^2$ exactly ($=2+\sqrt3$ for chain $A$,
$=1$ for chain $B$), confirming Lemma~\ref{lem:closure} independently;
and \emph{all} $66$ vertex pairs are classified exactly, reproducing
the $14$-edge sets used by the SAT certificates and confirming that all
$11$ charged slots are exactly at distance $r$ and are non-edges.  The
exact coordinates agree with the floating-point coordinate files
(\texttt{chain12.json}, \texttt{chain12b.json}) to $10^{-9}$.
\item \emph{Caps on exact edges.}  Proposition~\ref{prop:caps} is
re-proved on the exact edge lists by exhaustive enumeration: all
$\binom{11}{6}=462$ six-subsets of the slots are unsatisfiable as
simultaneous monochromatic sets under proper $6$-coloring, while a
five-subset is satisfiable; hence the maximum is exactly $5$.
\end{itemize}

\subsection{Proof of Theorem~\ref{thm:main}}

\begin{proof}[Proof of Theorem~\ref{thm:main}]
Suppose $c$ is a proper $6$-coloring of $G(\R^2;\{1,\dd\})$ and fix an
invariant mean $\Lambda$ (ZFC).  By Corollary~\ref{cor:v8},
$p_{\sqrt2}+p_{1+\sqrt3}\ge1$; by Corollary~\ref{cor:caps},
$p_{\sqrt2}+p_{1+\sqrt3}\le\tfrac{10}{11}$.  Contradiction.  Hence no
proper $6$-coloring exists and $\chid{\{1,\dd\}}\ge7$.
\end{proof}

\subsection{The framework caveat}\label{sec:caveat}

Two caveats separate Theorem~\ref{thm:main} from a classical
finite-witness bound such as de~Grey's $\chi(\R^2)\ge5$.  (A third ---
floating-point placement of the chain apexes --- was eliminated by the
exact-arithmetic certification of Section~\ref{sec:exact}: the V8 and
both chains are now verified exactly, end to end.)

\begin{enumerate}[label=(\roman*)]
\item \emph{Nonconstructive averaging.}  The densities $p_r$ exist by an
invariant mean on the amenable group $E(2)$, whose existence uses
Hahn--Banach.  This is a ZFC theorem, so Theorem~\ref{thm:main} rules
out \emph{arbitrary} (not merely measurable) $6$-colorings in ZFC.  But
in ZF+DC the mean need not exist, and Shelah--Soifer
\cite{ShelahSoifer03} showed chromatic numbers of distance graphs can
genuinely depend on the ambient set theory.  The choice-free content of
this section is exactly its three finite certificates.
\item \emph{No explicit finite witness yet.}  By the de~Bruijn--Erd\H{o}s
theorem a finite $7$-chromatic $\{1,\dd\}$-graph exists, but the argument
does not exhibit one.  Section~\ref{sec:blueprint} makes the
gap precise and reduces it to a combinatorial design problem;
Section~\ref{sec:verdict} reports computational evidence that the most
obvious design vocabulary does not suffice.
\end{enumerate}

\section{Toward a fully finite witness}\label{sec:blueprint}

\subsection{The counting blueprint}

The two halves of the proof are pointwise statements about finite
configurations; only the step from ``every placement'' to ``density''
uses the invariant mean.  A single finite graph would suffice if the
slots of many chain copies could be made to \emph{coincide} with the
diagonals of many V8 copies:

\begin{proposition}[Uniform-multiplicity blueprint]\label{prop:blueprint}
Suppose a finite point set $W$ contains $N$ V8 copies and a family
$\mathcal C$ of chain copies (of types $A$ and $B$) such that the
multiset of charged slots of $\mathcal C$ covers the $4N$ diagonal pairs
of the V8s, each with the same multiplicity $m\ge1$, and contains no
other pairs.  Then $G(W;\{1,\dd\})$ is not $6$-colorable.
\end{proposition}

\begin{proof}
Let $M$ be the number of monochromatic diagonals in a putative proper
$6$-coloring.  Summing Proposition~\ref{prop:v8} over the $N$ V8s gives
$M\ge2N$.  Each chain contributes $11$ slots of which at most $5$ are
monochromatic (Proposition~\ref{prop:caps}); summing over $\mathcal C$
and using the uniform multiplicity, $mM\le\tfrac5{11}\cdot 4Nm$, i.e.\
$M\le\tfrac{20}{11}N<2N$.  Contradiction.
\end{proof}

The design problem is thus to realize such an incidence structure with
actual plane geometry --- every slot at distance $\sqrt2$ or $1+\sqrt3$
being simultaneously a diagonal of some unit or outer square whose
V8 completion is present.  Exact uniform multiplicity is a strong
requirement; a defect-tolerant weighted version of
Proposition~\ref{prop:blueprint}, allowing some slots to be uncovered or
covered unevenly, appears feasible, but we have not written out a
complete proof and make no claim about it here.  A naive faithful closure
of the required incidences is F\o lner-type and very large, which
motivates:

\begin{question}[Gatekeeper question]\label{q:gatekeeper}
Does a finite $\{1,\dd\}$-graph exist containing two designated $r$-pairs
in a rational-angle relative position such that no proper $6$-coloring
makes both pairs monochromatic?  The natural target is ``not both
diagonals of one unit square monochromatic'' ($90^\circ$, self-supplied
by V8s).
\end{question}

A positive answer would allow far smaller finite witnesses than a naive
closure.  \emph{Computational observation} (not certified): we
searched closures of the $\zeta_{24}$ lattice (built from a $12$-fold
wheel of De~Neve-style gadgets about a hub, whose output ring is a
regular $12$-gon with side exactly $\sqrt2$ and $60^\circ$ chords
exactly $1+\sqrt3$), including neighborhoods of a designated unit square
grown for eight rounds inside a $24{,}881$-point universe, up to $1837$
points and $16{,}116$ edges.  In all of them, $6$-colorings making both
diagonals of the designated square monochromatic existed and were found
quickly.
Within this horizon we found no rational-angle cap gadget.

\subsection{Computational observations}\label{sec:verdict}

The following were obtained by exhaustive SAT enumeration and exact
symbolic computation.  The code and logs are not part of the certificate
bundle and the arguments have not been written up, so we report them as
observations that guided the search, not as theorems.

\begin{itemize}
\item \emph{Duplication and trees.}  Two copies of a configuration
sharing slots contributed at most one useful clause in all cases tested,
and tree-like slot incidence structures never forced a monochromatic
pair.
\item \emph{No short angle relations.}  No integer relation of weight
$\le7$ was found among $\alpha=\arccos\tfrac{2+\sqrt3}4$,
$\beta=\arccos\tfrac34$, $\gamma=\arccos\tfrac{2-\sqrt3}4$,
$\delta'=\arccos\tfrac{3\sqrt6-\sqrt2}8$, and $30^\circ$, so there are
no short ``odd spindle'' cycles among the relevant angles.
\item \emph{Two-V8 sharing.}  For two V8s sharing $4$ points, all
$121=11^2$ joint weight-$\ge2$ diagonal patterns were feasible, so
pattern-level reasoning over such pairs of V8s forces nothing.
\item \emph{Pattern sets.}  The sets of achievable monochromatic-slot
patterns were enumerated for the V8 (all patterns of weight $\ge2$, as in
Proposition~\ref{prop:v8}) and for both chains.  Restricting chain $A$ to
a $7$-slot subfamily gives cap $3/7$, the lowest cap ratio among all
$2047$ nonempty slot subsets.
\item \emph{Coincidences.}  Coincidence scans over the hinge families
($A$--$A$, $B$--$B$, and $A$--$B$ at $45^\circ$) found that two chain
placements share at most one charged slot, and the translation modules
governing periodic ladder assemblies had full $\Z$-rank in exact
arithmetic over $\Q(\sqrt3,\sqrt7)$ and the towers $K_A$, $K_B$ of
Section~\ref{sec:exact}.  This suggests that slot multiplicity cannot be
accumulated cheaply by crossing chains.
\item \emph{Linear-programming duals.}  A Farkas-dual linear program over
the enumerated pattern sets, with the coincidence constraints above, was
infeasible for the inventory consisting of V8s, chains and trimmed
chains, suggesting that these gadgets alone do not yield a finite
contradiction.
\end{itemize}

If these observations are correct, a finite witness needs further
ingredients, for example \emph{generic-position pigeonhole gadgets}
(V8-style pigeonholes operating in generic relative position, which is
also the direction of Question~\ref{q:gatekeeper}) or \emph{replacement
cap gadgets at new radii}, such as the star-polygon family of
Section~\ref{sec:k5caps}.

\section{The record six-chromatic graphs}\label{sec:record}

\subsection{A 26-vertex 6-chromatic $\{1,\dd\}$-graph}\label{sec:g26}

\begin{theorem}\label{thm:g26}
The $26$ points of Table~\ref{tab:g26} span $85$ pairs at distance
exactly $1$ ($48$ pairs) or exactly $\dd$ ($37$ pairs), and the graph on
these pairs has chromatic number $6$ and is vertex-critical.  Hence
$\chid{\{1,\dd\}}\ge6$ via a $26$-vertex witness --- to our knowledge
the smallest $6$-chromatic two-distance graph known for any ratio.
Certificates: \texttt{certificates/g26/}.
\end{theorem}

Each vertex is an integer $8$-vector $[a_0,\dots,a_7]$ encoding the point
\[
  \Bigl(\tfrac{a_0+a_1\sqrt3+a_2\sqrt7+a_3\sqrt{21}}{8},\
        \tfrac{a_4+a_5\sqrt3+a_6\sqrt7+a_7\sqrt{21}}{8}\Bigr)
  \in \Q(\sqrt3,\sqrt7)^2 .
\]
Distances are verified in exact integer arithmetic: writing a coordinate
difference as $(p,q,r,s)$ in the basis $(1,\sqrt3,\sqrt7,\sqrt{21})/8$,
its square contributes
\[
  \bigl(p^2+3q^2+7r^2+21s^2,\ 2pq+14rs,\ 2pr+6qs,\ 2ps+2qr\bigr)
\]
in the basis $(1,\sqrt3,\sqrt7,\sqrt{21})/64$; summing the $x$ and $y$
contributions, a pair is a unit edge iff the total is $(64,0,0,0)$ and a
$\dd$-edge iff it is $(128,64,0,0)$, i.e.\ squared distance
$2+\sqrt3=\dd^2$.

\begin{table}[ht]
\centering\small
\begin{tabular}{@{}rr@{}}
\toprule
\multicolumn{2}{c}{Vertices of the $26$-vertex graph
  ($[a_0,a_1,a_2,a_3\,|\,a_4,a_5,a_6,a_7]$)}\\
\midrule
$[-11,0,0,1\,|\,0,3,1,0]$    & $[-3,3,1,1\,|\,3,3,1,-1]$ \\
$[-8,-1,1,0\,|\,-1,0,0,-1]$  & $[-3,4,0,1\,|\,-4,3,1,0]$ \\
$[-8,0,0,0\,|\,0,0,0,0]$     & $[-3,4,0,1\,|\,4,3,1,0]$ \\
$[-8,3,-1,0\,|\,-3,0,0,-1]$  & $[-3,4,2,1\,|\,2,3,1,0]$ \\
$[-8,3,1,0\,|\,-5,0,0,-1]$   & $[-2,3,-1,0\,|\,-3,0,-2,-1]$ \\
$[-8,3,1,0\,|\,3,0,0,-1]$    & $[0,3,-1,0\,|\,-3,0,0,-1]$ \\
$[-7,0,0,1\,|\,0,-1,1,0]$    & $[0,3,1,0\,|\,-5,0,0,-1]$ \\
$[-7,4,0,1\,|\,-4,-1,1,0]$   & $[0,3,1,0\,|\,3,0,0,-1]$ \\
$[-6,0,0,0\,|\,0,0,2,0]$     & $[0,4,0,0\,|\,-4,0,0,0]$ \\
$[-6,1,1,0\,|\,-1,0,2,1]$    & $[0,7,1,0\,|\,-1,0,0,-1]$ \\
$[-6,3,-1,0\,|\,-3,0,2,-1]$  & $[1,4,0,1\,|\,-4,-1,1,0]$ \\
$[-6,4,0,0\,|\,-4,0,2,0]$    & \\
$[-5,3,1,1\,|\,3,3,-1,-1]$   & \\
$[-4,3,1,0\,|\,3,4,0,-1]$    & \\
$[-3,0,0,1\,|\,0,3,1,0]$     & \\
\bottomrule
\end{tabular}
\caption{Integer coordinate vectors of the $26$-vertex $6$-chromatic
$\{1,(\sqrt6+\sqrt2)/2\}$-graph.  Point encoded by
$\bigl((a_0+a_1\sqrt3+a_2\sqrt7+a_3\sqrt{21})/8,\
(a_4+a_5\sqrt3+a_6\sqrt7+a_7\sqrt{21})/8\bigr)$.}
\label{tab:g26}
\end{table}

\subsection{A 30-vertex 6-chromatic $\{1,\sqrt3\}$-graph}\label{sec:g30}

\begin{theorem}\label{thm:g30}
The $30$ points of Table~\ref{tab:g30} span $100$ pairs at distance
exactly $1$ ($72$) or exactly $\sqrt3$ ($28$), and the graph on these
pairs has chromatic number $6$ and is vertex-critical.  Hence
$\chid{\{1,\sqrt3\}}\ge6$ via a $30$-vertex witness (previously $33$
\cite{DeNeve25}).  Certificates: \texttt{certificates/g30/}.
\end{theorem}

Here each vertex is an integer $4$-vector $[a_0,a_1,a_2,a_3]$ encoding
\[
  \Bigl(\tfrac{a_0+a_1\sqrt5}{8},\ \tfrac{a_2\sqrt3+a_3\sqrt{15}}{8}\Bigr)
  \in\Q(\sqrt3,\sqrt5)^2,
  \qquad a_0\equiv a_1,\ a_2\equiv a_3 \pmod 2 .
\]
For a difference vector $(p,q,r,s)$,
$64\cdot\mathrm{dist}^2=p^2+5q^2+3r^2+15s^2+(2pq+6rs)\sqrt5$, so a unit
edge means exactly $(64,0)$ and a $\sqrt3$-edge exactly $(192,0)$.

\begin{table}[ht]
\centering\small
\begin{tabular}{@{}rrr@{}}
\toprule
\multicolumn{3}{c}{Vertices of the $30$-vertex graph
  ($[a_0,a_1,a_2,a_3]$)}\\
\midrule
$[-6,0,-2,0]$  & $[-1,-3,-1,1]$ & $[1,3,-1,1]$ \\
$[-5,-3,-1,1]$ & $[-1,-3,1,3]$  & $[1,3,3,1]$ \\
$[-5,3,-3,1]$  & $[-1,3,-1,3]$  & $[2,0,-2,0]$ \\
$[-4,0,-2,2]$  & $[-1,3,1,1]$   & $[2,0,0,2]$ \\
$[-4,0,0,0]$   & $[-1,3,3,3]$   & $[2,0,2,0]$ \\
$[-4,0,0,4]$   & $[0,0,-2,2]$   & $[3,-3,1,3]$ \\
$[-3,-3,-1,3]$ & $[0,0,0,0]$    & $[3,3,-1,3]$ \\
$[-3,-3,1,1]$  & $[0,0,2,2]$    & $[4,0,0,4]$ \\
$[-3,3,-1,1]$  & $[1,-3,-3,1]$  & \\
$[-2,0,-2,4]$  & $[1,-3,1,1]$   & \\
$[-2,0,0,2]$   & \\
$[-1,-3,-3,3]$ & \\
\bottomrule
\end{tabular}
\caption{Integer coordinate vectors of the $30$-vertex $6$-chromatic
$\{1,\sqrt3\}$-graph.  Point encoded by
$\bigl((a_0+a_1\sqrt5)/8,\ (a_2\sqrt3+a_3\sqrt{15})/8\bigr)$.}
\label{tab:g30}
\end{table}

\subsection{Method: SAT-guided iterated local search over module
closures}\label{sec:method}

Both graphs were found by the same pipeline, seeded from the published
coordinate tables of De~Neve et al.\ \cite{DeNeve25} (their $33$-vertex
$\{1,\sqrt3\}$-graph and $117$-vertex $\{1,\dd\}$-graph).

\begin{itemize}
\item \emph{Universe.}  A bounded-coefficient slice of the $\Z$-module
closure of the seed coordinates: all points $p+q-r$ over the current
point set staying in the same coset of the coordinate module (for
$\{1,\sqrt3\}$: $\{((a_0+a_1\sqrt5)/8,(a_2\sqrt3+a_3\sqrt{15})/8):
a_i\in\Z,\ a_0\equiv a_1,\ a_2\equiv a_3\ (\mathrm{mod}\ 2)\}$,
$\approx10^4$ points within radius~$3.4$).  The closure contains
coefficient patterns the seed constructions never generate; the
$26$-vertex graph uses several such points.
\item \emph{Oracle.}  An incremental CaDiCaL instance with per-vertex
selector assumptions answers ``is this subset still $6$-chromatic?'' in
microseconds.
\item \emph{Search.}  Randomized greedy vertex deletion descends to
critical subgraphs; perturbation kicks (adjacency-weighted candidate
insertion, remove-$k$, re-descend) escape local minima.  The $30$-vertex
minimum was reached within $\sim10^4$ oracle calls ($\approx1$ minute on
an Apple M4~Pro); it shares $23$ vertices with the $33$-vertex seed,
omits $10$, and adds $7$ new.  In the $\{1,\dd\}$ search the descent
first stopped at $29$ vertices because the clique used for symmetry
breaking had been exempted from deletion; a final pass without this
exemption reduced the graph to $26$ vertices.
\end{itemize}

\subsection{Verification recipe}\label{sec:verify}

Every claim in Theorems~\ref{thm:g26} and~\ref{thm:g30} is re-derivable
from the integer vector tables alone; the files are in
\texttt{certificates/g26/} and \texttt{certificates/g30/}.

\begin{enumerate}
\item \emph{Exact edges.}  Recompute all pairwise squared distances in
integer arithmetic as above; no floating point enters any certificate.
\item \emph{Non-$5$-colorability.}  The direct $130$-- resp.\
$150$-variable CNF (one variable per vertex-color pair, no
symmetry-breaking assumptions; \texttt{record26.cnf},
\texttt{record30.cnf}) is UNSAT, with DRAT traces
\texttt{record26.drat} ($18{,}518$ lines) and \texttt{record30.drat}
(compressed as \texttt{record30.drat.xz}) checked by \texttt{drat-trim}
(\texttt{s VERIFIED}).
\item \emph{$6$-colorability and criticality.}  Explicit $6$-colorings
exist, and all single-vertex-deleted subgraphs are $5$-colorable; both
are found by SAT in the verifiers.
\item \emph{Standalone verifiers.}  \texttt{verify117best.py
deneve117\_min.json} and \texttt{verify30.py} ($\sim60$ lines each;
dependency: \texttt{python-sat}) recompute the edge sets exactly from the
integer vectors and re-solve non-$5$-colorability (for the $26$-vertex
graph with two different SAT solvers), $6$-colorability and
vertex-criticality from scratch in seconds.
\end{enumerate}

\begin{remark}
Whether a $25$-vertex (or smaller) $6$-chromatic two-distance graph
exists for any ratio is open.  Note that the bound $v_6\ge42$ for
unit-distance graphs \cite{deGreyParts22} does not apply to two-distance
graphs, where only trivial lower bounds are known for $\{1,\sqrt3\}$.
\end{remark}

\section{Unit-distance results}\label{sec:ud}

\subsection{$e_5\le2278$}\label{sec:e5}

The smallest known $5$-chromatic unit-distance graph has $509$ vertices
and $2442$ edges \cite{Parts20b,deGreyParts22}.  It was known not to be
edge-critical: de~Grey and Parts removed $36$ edges non-exhaustively,
giving $e_5\le2406$, and noted room for improvement
\cite{deGreyParts22}.

\begin{theorem}\label{thm:e5}
The $509$-vertex graph contains a spanning subgraph with $2278$ edges
($164$ removed) that is not $4$-colorable.  Hence $e_5\le2278$.
\end{theorem}

The subgraph is listed in \texttt{certificates/e5-2278/edge2278\_graph.json}
(an edge list on the vertex numbering $0,\dots,508$ of the $509$-vertex
graph).  Its non-$4$-colorability is certified by
\texttt{edge2278.cnf}/\texttt{edge2278.drat} (compressed with
\texttt{xz}), checked by \texttt{drat-trim} in about ten seconds.  The
CNF has one at-least-one-color clause per vertex, one clause per edge
and color, and three unit clauses fixing the colors of the three
vertices of one triangle of the subgraph, which is a valid symmetry
breaking.  The bundle also contains the exact algebraic coordinates of the $509$
vertices (\texttt{parts509\_coordinates.vtx}, from~\cite{Parts20b}) and a
script \texttt{check\_unit.py} which confirms, at $60$-digit precision, that
these points have exactly $2442$ unit-distance pairs, with no near-unit pair,
and that the $2278$ listed edges are among them.

The trim used the same incremental selector-assumption technique as
Section~\ref{sec:method}, over edges instead of vertices: repeatedly
attempt to delete a random edge, keeping the deletion iff the remainder
is still not $4$-colorable, until no edge is removable.
\emph{Computational observation} (not part of the certificate bundle):
the final subgraph is edge-critical, i.e.\ removing any one of its $2278$
edges leaves a $4$-colorable graph; a second, independently seeded run
terminated at $2279$ edges, so the greedy endpoint is order-dependent
and the true $e_5$ may be lower still.

\subsection{Forced structure in $5$-colorings of the plane}\label{sec:forced}

The transfer principle of Section~\ref{sec:framework} applies equally
with $k=5$ and $D=\{1\}$: a non-$5$-colorable $\{1,\dd'\}$-graph whose
$m$ $\dd'$-edges cannot all be avoided forces monochromatic pairs at
distance $\dd'$ in \emph{every} proper unit-distance $5$-coloring of the
plane (cf.\ Gwyn--Stavrianos \cite{GwynStavrianos20}).  Since
$\chi(\R^2)\ge6$~\cite{OAI}, the following corollary is vacuous; we
record it because its finite content --- the ratio $t/m$ of forced to
designated pairs in an explicit point set --- was the original
motivation for the search for small graphs, and because it is the
quantity in the win condition of Section~\ref{sec:k5caps}.

\begin{corollary}\label{cor:forced}
In every proper $5$-coloring of $G(\R^2;\{1\})$, with $p_r$ the invariant
densities of Definition~\ref{def:pr}:
\[
  p_{\sqrt3}\ \ge\ \tfrac1{28},
  \qquad
  p_{\dd}\ \ge\ \tfrac1{37},
  \qquad
  p_{1/\sqrt3}\ \ge\ \tfrac1{67}.
\]
\end{corollary}

\begin{proof}
Consider the $30$-point set of Theorem~\ref{thm:g30}.  Any $5$-coloring
of the plane proper at distance $1$ restricts to a coloring of the $30$
points proper on the $72$ unit edges; since the full graph is not
$5$-colorable (\texttt{g30/}), at least one of the $28$ $\sqrt3$-pairs is
monochromatic.  By Lemma~\ref{lem:transfer}, $28\,p_{\sqrt3}\ge1$.  The
same argument on the $26$-point set (Theorem~\ref{thm:g26},
\texttt{g26/}; $37$ $\dd$-pairs) gives $37\,p_{\dd}\ge1$.

For the third bound, \texttt{certificates/forced-density/} contains a set
$S$ of $67$ of the $72$ unit pairs of the same $30$-point set such that
the graph formed by all $28$ $\sqrt3$-pairs together with $S$ is not
$5$-colorable (\texttt{record67.cnf}/\texttt{record67.drat}, checked by
\texttt{drat-trim}; \texttt{verify67.py} also recomputes all distances
exactly and checks that $S$ is minimal).  Scale the point set by
$1/\sqrt3$: the $\sqrt3$-pairs become unit pairs and the pairs of $S$
become pairs at distance $1/\sqrt3$.  For every placement, a coloring of
the plane proper at distance $1$ is proper on the former $\sqrt3$-pairs,
so at least one pair of $S$ is monochromatic, and
Lemma~\ref{lem:transfer} gives $67\,p_{1/\sqrt3}\ge1$.
\end{proof}

For comparison, Parts' $397$-vertex $\{1,2\}$-graph \cite{Parts20a}
implicitly forces $p_2$ to be of order $1/550$ in the same sense.
Conversely $p_{\dd'}\le\tfrac12$ for all $\dd'\ge\tfrac12$ by a
color-count-free argument: if $|xa|=|xb|=\dd'$ and $|ab|=1$, at most one
of the pairs $(x,a)$, $(x,b)$ is monochromatic.  So any construction with
$t/m>\tfrac12$ would have proved $\chi(\R^2)\ge6$, and
Section~\ref{sec:k5caps} lowers this threshold to $\tfrac13$ at the
distance $1/\sqrt3$.

\begin{remark}[Density ledger; computational observation]
\label{rem:ledger}
We also assembled about $143{,}000$ linear inequalities among
monochromatic-pair densities, each obtained from a small unit-distance
configuration as in Lemma~\ref{lem:transfer} with MaxSAT-computed
bounds, and solved the resulting linear program in floating point.  It
returned the bound $p_{\sqrt3}\ge1/12$ for invariant proper
$5$-colorings.  The inequalities were not individually certified and the
linear program was not re-solved in exact arithmetic, so we do not claim
this bound (which is in any case vacuous now).
\end{remark}

\begin{remark}[A single-distance sweep; computational observation]
\label{rem:wall}
For Heule's $803$-vertex $5$-chromatic unit-distance graph (April
2018, from the Polymath16 minimization effort
\cite{Polymath16,Heule18}), we tested, for each of the $428$ distance classes with at
least $40$ point pairs, whether some proper $5$-coloring avoids
monochromatic pairs in that class entirely.  The answer was yes for
every class, so no single distance class is forced monochromatic in
this graph.  All classes were resolved within seconds except the class
of distance $2$, which took $3.4$ hours of solver time before a witness
coloring was found.
\end{remark}

\section{Strict caps for five colors and a win condition}\label{sec:k5caps}

\emph{(This section was written before~\cite{OAI} proved
$\chi(\R^2)\ge6$, which makes the win condition of
Theorem~\ref{thm:win} moot.  We keep it for the cap statements, which
remain valid.)}

Corollary~\ref{cor:caps} caps a density strictly below $\tfrac12$ for
$6$-colorings.  This section develops the analogous theory on the
unit-distance side: caps valid for colorings with \emph{any} number of
colors, a bipartiteness obstruction explaining where such pairwise caps
cannot exist, and the resulting win condition for $\chi(\R^2)\ge6$.
Throughout, ``proper'' means proper at distance $1$, and $p_r$ is the
invariant density of Definition~\ref{def:pr} for such a coloring.

\subsection{Color-count-free finite caps}

\begin{proposition}[Triangle and circumcenter]\label{prop:cap13}
Let $Q$ consist of a unit equilateral triangle $v_1v_2v_3$ and its
circumcenter $o$, with $|o-v_i|=1/\sqrt3$.  In every proper coloring
(any number of colors), at most one of the three pairs $(o,v_i)$ is
monochromatic.  Hence $p_{1/\sqrt3}\le\tfrac13$ for every proper
coloring of the plane.
\end{proposition}

\begin{proof}
The $v_i$ are pairwise at distance $1$, so they receive three distinct
colors; $o$ can agree with at most one of them.  Transfer
(Lemma~\ref{lem:transfer}) gives $3\,p_{1/\sqrt3}\le1$.
\end{proof}

\begin{proposition}[Unit pentagon]\label{prop:cap25}
Let $\tau=(1+\sqrt5)/2$.  The regular pentagon with unit sides has
diagonals of length $\tau$; in every proper coloring, at most two of
the five diagonals are monochromatic, so $p_{\tau}\le\tfrac25$.
Dually, the regular pentagon with unit diagonals has sides of length
$1/\tau$ and caps $p_{1/\tau}\le\tfrac25$.
\end{proposition}

\begin{proof}
Label the vertices $v_0,\dots,v_4$ cyclically.  If two diagonals
sharing a vertex, say $(v_0,v_2)$ and $(v_2,v_4)$, were both
monochromatic, then $c(v_0)=c(v_4)$, contradicting the unit side
$(v_4,v_0)$.  So monochromatic diagonals form an independent set in the
diagonal $5$-cycle $v_0v_2v_4v_1v_3$, of size at most $2$.  The dual
statement is identical with sides and diagonals exchanged.
\end{proof}

\begin{proposition}[Odd stars and flowers]\label{prop:capstar}
For odd $n$ and $1\le a<n/2$ with $\gcd(a,n)=1$: place a regular
$n$-gon so that its step-$2a$ chords have length $1$; its step-$a$
chords then have length $r=1/(2\cos(a\pi/n))$, and in every proper
coloring at most $(n-1)/2$ of the $n$ step-$a$ chords are
monochromatic: $p_r\le\tfrac{n-1}{2n}<\tfrac12$.  Similarly the
``flower'' (regular $m$-gon with unit sides plus its center, $m$ odd)
caps $p_r\le\tfrac{m-1}{2m}$ at the circumradius $r=1/(2\sin(\pi/m))$;
$m=3$ recovers Proposition~\ref{prop:cap13}.  The capped radii
$1/(2\cos(a\pi/n))$ are dense in $(\tfrac12,\infty)$.
\end{proposition}

\begin{proof}
If the circumradius is $R$, the step-$k$ chords have length
$2R\sin(k\pi/n)$; with $2R\sin(2a\pi/n)=1$ the step-$a$ chords have
length $\sin(a\pi/n)/\sin(2a\pi/n)=1/(2\cos(a\pi/n))$.
Two step-$a$ chords sharing a vertex, $(v_i,v_{i+a})$ and
$(v_{i+a},v_{i+2a})$, cannot both be monochromatic: that would force
$c(v_i)=c(v_{i+2a})$ across the unit step-$2a$ chord.  So monochromatic
step-$a$ chords are independent in the step-$a$ cycle, which is a
single \emph{odd} cycle since $\gcd(a,n)=1$ and $n$ is odd; an odd
$n$-cycle has independence number $(n-1)/2$, and
Lemma~\ref{lem:transfer} gives the cap.  For the flower, consecutive
spokes $(o,v_i)$, $(o,v_{i+1})$ cannot both be monochromatic because
$|v_iv_{i+1}|=1$, and the $m$ spokes form an odd $m$-cycle under this
relation.  Finally, the fractions $a/n$ with $n$ odd, $\gcd(a,n)=1$ and
$a<n/2$ are dense in $(0,\tfrac12)$, so the radii $1/(2\cos(a\pi/n))$
are dense in $(\tfrac12,\infty)$.
\end{proof}

These caps hold for arbitrary proper colorings, with no bound on the
number of colors; they are elementary, and we do not claim priority for
them.

\subsection{A bipartiteness obstruction at $\sqrt3$ and $2$}

Why do the caps above live at $1/\sqrt3$, $\tau$, $1/(2\cos(a\pi/n))$
and not at $\sqrt3$ and $2$, the second distances of the
$\{1,\sqrt3\}$- and $\{1,2\}$-constructions?  All
of them use the same mechanism.  For $r\ge\tfrac12$, call an unordered
pair of points at distance $r$ an \emph{$r$-slot}, and say that two
$r$-slots $\{x,a\}$, $\{x,b\}$ sharing an endpoint \emph{conflict} if
$|a-b|=1$.  Two conflicting slots cannot both be monochromatic in a
proper coloring, since that would give $c(a)=c(x)=c(b)$.  Hence, for a
finite family $F$ of $r$-slots with conflict graph $\Gamma(F)$ (vertices
the slots, edges the conflicting pairs), the monochromatic slots of any
proper coloring form an independent set of $\Gamma(F)$, and
Lemma~\ref{lem:transfer} gives the \emph{pairwise cap}
\[
  p_r\ \le\ \frac{\alpha(\Gamma(F))}{|F|}.
\]
Propositions~\ref{prop:cap13}--\ref{prop:capstar} are all of this form,
with $\Gamma(F)$ an odd cycle.  The angle at the shared endpoint of two
conflicting slots is $\varphi_1(r)=2\arcsin\frac1{2r}$, with
$\cos\varphi_1=1-\tfrac1{2r^2}$ (at $r=1/\sqrt3$:
$\cos\varphi_1=-\tfrac12$, $\varphi_1=2\pi/3$; at $r=\tau$:
$\cos\varphi_1=\tau/2$, $\varphi_1=\pi/5$).

\begin{proposition}[Pairwise caps at $\sqrt3$ and $2$]
\label{prop:bipartite}
Let $r\ge\tfrac12$ and suppose $\varphi_1(r)/\pi$ is irrational.  Then
$\Gamma(F)$ is bipartite for every finite family $F$ of $r$-slots, so
$\alpha(\Gamma(F))\ge|F|/2$ and every pairwise cap at distance $r$ is
at least $\tfrac12$.  This applies to $r=\sqrt3$ and $r=2$, where
$\cos\varphi_1=\tfrac56$ and $\tfrac78$ respectively.
\end{proposition}

\begin{proof}
Assign to each slot $\{x,a\}$ the direction $\theta(\{x,a\})\in\R/\pi\Z$
of the line through $x$ and $a$.  If $\{x,a\}$ and $\{x,b\}$ conflict,
the triangle $xab$ has sides $r,r,1$, so the angle at $x$ is $\varphi_1$
and $\theta(\{x,b\})=\theta(\{x,a\})\pm\varphi_1$ in $\R/\pi\Z$.  Along a
closed walk of length $L$ in $\Gamma(F)$ we therefore get
$(\varepsilon_1+\dots+\varepsilon_L)\varphi_1\in\pi\Z$ with
$\varepsilon_i=\pm1$.  Since $\varphi_1/\pi$ is irrational,
$\varepsilon_1+\dots+\varepsilon_L=0$, so $L$ is even.  Hence
$\Gamma(F)$ has no odd cycle and is bipartite, and the larger side of a
bipartition is an independent set with at least $|F|/2$ slots.  For
$r=\sqrt3$ and $r=2$, $\cos\varphi_1=1-\tfrac1{2r^2}$ is rational and
not in $\{0,\pm\tfrac12,\pm1\}$, so by Niven's theorem \cite{Niven56}
$\varphi_1/\pi$ is irrational.
\end{proof}

Thus at $\sqrt3$ and $2$ a cap below $\tfrac12$ needs constraints
beyond pairwise slot conflicts.  \emph{Computational observation} (not
certified): SAT computations on scaffolds of up to $45$ points and
exhaustive scans of chain-like configurations at $r=\sqrt3$ and $r=2$
never gave a cap below $\tfrac12$, while the same scans at $r=1/\sqrt3$
found chains with cap $3/8$.

\subsection{The win condition}

Before~\cite{OAI}, the caps converted record-graph engineering into a
quantitative attack on $\chi(\R^2)\ge6$.

\begin{theorem}[Win condition]\label{thm:win}
Let $G$ be a finite $\{1,\sqrt3\}$-graph with $m$ unit edges such that
every $5$-coloring of its vertices that is proper on the
$\sqrt3$-edges makes at least $t$ unit edges monochromatic (any
$6$-chromatic $\{1,\sqrt3\}$-graph has $t\ge1$).  If
$t/m>\tfrac13$, then $\chi(\R^2)\ge6$.
\end{theorem}

\begin{proof}
Suppose $c$ is a proper $5$-coloring of $G(\R^2;\{1\})$.  Scale $G$ by
$1/\sqrt3$: its $\sqrt3$-edges become unit pairs and its unit edges
become pairs at distance $1/\sqrt3$.  For every isometric placement of
the scaled vertex set, $c$ restricted to it is a $5$-coloring proper on
the former $\sqrt3$-edges, so by hypothesis at least $t$ of the $m$
former unit edges are monochromatic.  Lemma~\ref{lem:transfer} gives
$p_{1/\sqrt3}\ge t/m>\tfrac13$, contradicting
Proposition~\ref{prop:cap13}.  Hence no proper $5$-coloring exists.
\end{proof}

The best certified value of the ratio is $t/m=1/67$
(\texttt{certificates/forced-density/}; see the proof of
Corollary~\ref{cor:forced}): the $30$-vertex graph of
Theorem~\ref{thm:g30} with its $28$ $\sqrt3$-edges and a minimal set of
$67$ of its $72$ unit edges is already not $5$-colorable, improving the
ratio $1/72$ of the full graph.  The golden-ratio variant uses
Proposition~\ref{prop:cap25}: a $\{1,\tau\}$-graph forcing unit-edge
monochromatic density above $\tfrac25$ under $\tau$-proper
$5$-colorings would also have proved $\chi(\R^2)\ge6$.  The gap between
$1/67$ and $1/3$ is large, and the question is now moot.

\section{An extended six-coloring continuum}\label{sec:continuum}

Following Mundinger, Pokutta, Spiegel and Zimmer~\cite{MPSZ}, let $X_6$ be the set of $d>0$ for which the plane has a coloring with six colors in which five color classes contain no two points at distance $1$ and the sixth contains no two points at distance $d$. They showed $X_6\supseteq[0.354,0.657]$ (to three decimals). Reaching $d=1$ would prove $\chi(\R^2)\le6$.

\begin{theorem}\label{thm:continuum}
$X_6\supseteq[0.354,\ 0.667504532081]$.
\end{theorem}

The new part is $[657/1000,\ 667504532081/10^{12}]$; the rest is
\cite{MPSZ}.  It is covered by $59$ periodic tilings with rational
vertex coordinates, each with eight polygonal tiles per period (five
tile types for the five unit-avoiding colors and three for the sixth
color), taken from a $33$-parameter family without mirror symmetry.  The
certificates are in \texttt{certificates/six-coloring-continuum/}.  For
each tiling, the exact rational-arithmetic verifier
\texttt{rigorous\_verify.py} checks that all tiles are simple polygons,
that periodic interiors do not overlap, that the tile areas sum exactly
to the area of the lattice cell, that each of the five unit-avoiding tile
types has diameter $<1$ with distinct periodic copies more than $1$
apart, and that the components of the sixth color have diameter below,
and mutual separation above, an explicit rational interval of distances
$d$.  The sixth color then avoids every $d$ in that open interval (with
any deterministic assignment of boundary points).  The script
\texttt{verify\_continuum.py} checks that the $59$ intervals overlap and
cover the stated range, and \texttt{verify\_all.py} regenerates all $59$
certificates from the stored parameter vectors before running the
coverage check.

\begin{remark}[Numerical observations]
The following are floating-point observations and are not part of the
proof of Theorem~\ref{thm:continuum}.  The numerical optimum of the
$33$-parameter family is $d=0.66750453208133\ldots$; there $23$
constraints are active with positive multipliers and the reduced Hessian
is numerically negative definite, indicating a strict local maximum for
this tiling topology.  Period doubling, $2\times2$, $3\times1$,
$1\times3$ and $3\times2$ supercells, and all period-two relabelings of
the five unit colors returned the same value.
\end{remark}

\section{Discussion and open problems}\label{sec:open}

\paragraph{Open problems.}

\begin{problem}
Is there a $25$-vertex $6$-chromatic two-distance graph?  The
$26$-vertex graph is vertex-critical and further local search did not
reduce it; no nontrivial lower bound on the minimum order is known for
two-distance graphs.
\end{problem}

\begin{problem}[Finite $7$-chromatic witness]
Realize the incidence structure of Proposition~\ref{prop:blueprint} by
an actual finite $\{1,\dd\}$-point set, or answer
Question~\ref{q:gatekeeper}.  A finite $\{1,\dd\}$-graph $W$ containing
two points at distance $s\ge\tfrac12$ that receive the same color in
every proper $6$-coloring would yield a finite non-$6$-colorable graph
on fewer than $2|W|$ vertices: take two copies of $W$ sharing one of the
two points, rotated so that the two partner points are at distance $1$.
\end{problem}

\begin{problem}
Determine $e_5$.  The greedy endpoint $2278$ is order-dependent
($2279$ on a second run), and we know no nontrivial lower bound
technique for the number of edges of $5$-chromatic unit-distance
graphs.
\end{problem}

\appendix

\section{Artifacts and reproducibility}\label{app:artifacts}

All certificates are in the folder \texttt{certificates/} next to this
manuscript; its \texttt{README.md} lists the check for each result, and
all checks passed on 14 October 2026.  Trusting an UNSAT claim requires
trusting only the DRAT checker \texttt{drat-trim} \cite{drattrim}
(available from \url{https://github.com/marijnheule/drat-trim}) and the
short scripts that build or read the CNF files; trusting the SAT
(colorability) claims requires only checking explicit colorings; the
exact-arithmetic chain certificates additionally rest on
\texttt{sympy}'s multivariate polynomial reduction and
minimal-polynomial routines; the tiling certificates use only Python's
exact \texttt{Fraction} arithmetic.  The searches and the computational
observations of Sections~\ref{sec:blueprint}, \ref{sec:ud}
and~\ref{sec:k5caps} are not part of the bundle.  Files ending in
\texttt{.xz} must be decompressed first (\texttt{xz -dk}).

\begin{center}
\small
\begin{tabular}{@{}p{4.4cm}lp{7.0cm}@{}}
\toprule
File(s) & Size & Content \\
\midrule
\multicolumn{3}{@{}l}{\texttt{seven-colors/} (Theorem~\ref{thm:main})}\\
\texttt{v8cert.cnf/.drat} & 2.7\,KB / 262\,KB &
proper $6$-coloring of the V8 with $\le1$ monochromatic diagonal:
UNSAT (Proposition~\ref{prop:v8}); $84$ vars, $267$ clauses. \\
\texttt{capA.cnf/.drat} & 7.2\,KB / 21\,KB &
proper $6$-coloring of chain $A$ ($r=\sqrt2$) with $\ge6$ of $11$ slots
monochromatic: UNSAT (Proposition~\ref{prop:caps}); $188$ vars, $646$
clauses. \\
\texttt{capB.cnf/.drat} & 7.2\,KB / 21\,KB &
same for chain $B$ ($r=1+\sqrt3$). \\
\texttt{chain12.json, chain12b.json, v8.json} & $<$1\,KB each &
floating-point coordinates and designated pairs of the three
configurations. \\
\texttt{verify\_exact.py} & $\sim$250 lines &
exact symbolic certification of chains $A$ and $B$
(Section~\ref{sec:exact}): tower-degree certificates via
primitive-element minimal polynomials, all $66$ pairwise distances per
chain resolved in $K_A$ resp.\ $K_B$, edge and slot sets confirmed
against the float files, caps re-proved by exhaustive subset
enumeration (dependencies: \texttt{sympy}, \texttt{python-sat}). \\
\midrule
\multicolumn{3}{@{}l}{\texttt{g26/} (Theorem~\ref{thm:g26}; $p_{\dd}\ge1/37$)}\\
\texttt{deneve117\_min.json} & $<$1\,KB &
the $26$ integer vectors (Table~\ref{tab:g26}). \\
\texttt{record26.cnf/.drat} & 4.9\,KB / 884\,KB &
$5$-coloring of the $26$-vertex graph: UNSAT. \\
\texttt{verify117best.py} & $\sim$60 lines &
exact edges, non-$5$-colorability (two solvers), $6$-colorability,
vertex-criticality. \\
\midrule
\multicolumn{3}{@{}l}{\texttt{g30/} (Theorem~\ref{thm:g30}; $p_{\sqrt3}\ge1/28$)}\\
\texttt{record30.cnf}, \texttt{record30.drat.xz} & 5.8\,KB / 1.3\,MB &
$5$-coloring of the $30$-vertex graph: UNSAT. \\
\texttt{verify30.py} & $\sim$70 lines &
exact edges, non-$5$-colorability, $6$-colorability,
vertex-criticality. \\
\midrule
\multicolumn{3}{@{}l}{\texttt{forced-density/} ($p_{1/\sqrt3}\ge1/67$)}\\
\texttt{record67\_graph.json} & 1\,KB &
the $30$ integer vectors and the $67$ designated unit pairs. \\
\texttt{record67.cnf/.drat} & 8.8\,KB / 776\,KB &
$5$-coloring proper on all $\sqrt3$-pairs and the $67$ unit pairs:
UNSAT. \\
\texttt{verify67.py} & $\sim$120 lines &
exact distances, non-$5$-colorability, regeneration of the CNF and DRAT
files, \texttt{drat-trim} check, minimality of the $67$ pairs. \\
\midrule
\multicolumn{3}{@{}l}{\texttt{e5-2278/} (Theorem~\ref{thm:e5})}\\
\texttt{edge2278\_graph.json} & 26\,KB &
edge list of the $2278$-edge subgraph of the $509$-vertex graph. \\
\texttt{edge2278.cnf.xz/.drat.xz} & 123\,KB / 70\,MB &
$4$-coloring of the subgraph (one triangle's colors fixed): UNSAT
(sizes uncompressed). \\
\midrule
\multicolumn{3}{@{}l}{\texttt{six-coloring-continuum/} (Theorem~\ref{thm:continuum})}\\
\texttt{certificate\_*.json}, \texttt{asym\_*.npy} & --- &
the $59$ exact tiling certificates and their parameter vectors. \\
\texttt{rigorous\_verify.py}, \texttt{verify\_all.py},
\texttt{verify\_continuum.py} & --- &
exact check of one tiling; regeneration of all $59$ certificates;
interval-coverage check. \\
\bottomrule
\end{tabular}
\end{center}

\noindent
Reproduction: \texttt{pip install python-sat sympy numpy}; then, in the
respective folders, \texttt{python3 verify\_exact.py},
\texttt{python3 verify117best.py deneve117\_min.json},
\texttt{python3 verify30.py},
\texttt{python3 verify67.py record67\_graph.json} (which looks for
\texttt{drat-trim} on the \texttt{PATH}, in the environment variable
\texttt{DRAT\_TRIM}, or as a second argument, and writes its CNF and DRAT
files to the current directory), and
\texttt{python3 verify\_continuum.py}; for any UNSAT certificate,
\texttt{drat-trim X.cnf X.drat} and confirm \texttt{s VERIFIED}.  All
searches ran on a single laptop (Apple M4~Pro); total compute for the
record graphs was minutes and for the edge trim hours.

\section*{Acknowledgments}
We thank Jaan Parts, whose 2019 sketch \cite{Parts19} contains the
skeleton of Theorem~\ref{thm:main}, and De~Neve, Vanden Kerchove, Colle,
Tavernier, and Pickavet, whose published coordinate tables
\cite{DeNeve25} seeded the searches of Section~\ref{sec:record}.
The searches, certificates and verification code were developed with the
assistance of Claude (Anthropic).

\begin{thebibliography}{99}

\bibitem{Agoston21}
P.~\'Agoston,
\emph{Probabilistic formulation of the Hadwiger--Nelson problem},
arXiv:2112.07665 (2021).

\bibitem{BassMono}
H.~Bass,
\emph{Monochromatic distances in finite colorings of the plane},
preprint (October 2026).

\bibitem{BassR3}
H.~Bass,
\emph{A machine-checked proof that three-space is not six-colorable},
preprint (October 2026).

\bibitem{BassRamsey}
H.~Bass,
\emph{Euclidean Ramsey theorems for arbitrary colourings via spectral
rigidity}, preprint (October 2026).

\bibitem{Bukh}
B.~Bukh,
\emph{Measurable sets with excluded distances},
Geom.\ Funct.\ Anal.\ \textbf{18} (2008), 668--697; arXiv:math/0703856.

\bibitem{Crocker}
C.~Crocker,
\emph{A conditional transfer for colourings of $\R^n$, with the
conditional consequence $\chi(\R^3)\ge7$}, Zenodo (10 October 2026),
doi:10.5281/zenodo.23273952.

\bibitem{deGrey18}
A.~D.~N.~J. de~Grey,
\emph{The chromatic number of the plane is at least 5},
Geombinatorics \textbf{28}(1) (2018), 18--31; arXiv:1804.02385.

\bibitem{deGreyParts22}
A.~D.~N.~J. de~Grey and J.~Parts,
\emph{On lower bounds of the order of $k$-chromatic unit distance
graphs}, Geombinatorics \textbf{32}(2) (2022), 72--74; arXiv:2303.14714.

\bibitem{DeNeve25}
J.~De~Neve, F.~Vanden~Kerchove, D.~Colle, W.~Tavernier, and
M.~Pickavet,
\emph{On the chromatic number of the plane with two forbidden
distances}, Amer.\ Math.\ Monthly \textbf{132}(10) (2025), 1007--1022;
doi:10.1080/00029890.2025.2559554.

\bibitem{ExooIsmailescu18}
G.~Exoo and D.~Ismailescu,
\emph{The chromatic number of the plane is at least 5: a new proof},
Discrete Comput.\ Geom.\ \textbf{64} (2020), 216--226; arXiv:1805.00157.

\bibitem{ExooIsmailescu20}
G.~Exoo and D.~Ismailescu,
\emph{A 6-chromatic two-distance graph in the plane},
Geombinatorics \textbf{29}(3) (2020); arXiv:1909.13177.

\bibitem{GwynStavrianos20}
H.~Gwyn and J.~Stavrianos,
\emph{A finite graph approach to the probabilistic Hadwiger--Nelson
problem}, arXiv:2008.07987 (2020).

\bibitem{Heule18}
M.~J.~H. Heule,
\emph{Computing small unit-distance graphs with chromatic number 5},
Geombinatorics \textbf{28}(1) (2018), 32--50; arXiv:1805.12181.

\bibitem{MPSZ}
K.~Mundinger, S.~Pokutta, C.~Spiegel, and M.~Zimmer,
\emph{Extending the continuum of six-colorings}, arXiv:2404.05509
(2024).

\bibitem{Niven56}
I.~Niven,
\emph{Irrational Numbers},
Carus Mathematical Monographs 11, Math.\ Assoc.\ America, 1956.

\bibitem{OAI}
OpenAI,
\emph{The Euclidean plane is not five-colorable}, preprint (2026),
\url{https://github.com/openai/math}, with Lean formalization.

\bibitem{OTH08}
J.~Owings, M.~Tetiva, and M.~Huddleston,
\emph{Coloring the plane}, Amer.\ Math.\ Monthly \textbf{115}(2) (2008),
170--172.

\bibitem{PalvolgyiAgoston19}
D.~P\'alv\"olgyi and P.~\'Agoston, comments in Polymath16, fourteenth
thread, comment 24460 (21 October 2019), on D.~G.~Mixon's blog
\emph{Short, Fat Matrices},
\url{https://dustingmixon.wordpress.com/2019/08/05/polymath16-fourteenth-thread-automated-graph-minimization/}.

\bibitem{Parts19}
J.~Parts, comment 24686 (24 December 2019) on \emph{Polymath16,
fifteenth thread: writing the paper and chasing down loose ends},
D.~G.~Mixon's blog \emph{Short, Fat Matrices},
\url{https://dustingmixon.wordpress.com/2019/12/12/polymath16-fifteenth-thread-writing-the-paper-and-chasing-down-loose-ends/#comment-24686}.

\bibitem{Parts20a}
J.~Parts,
\emph{A small 6-chromatic two-distance graph in the plane},
Geombinatorics \textbf{29}(3) (2020), 111--115; arXiv:2010.12656.

\bibitem{Parts20b}
J.~Parts,
\emph{Graph minimization, focusing on the example of 5-chromatic
unit-distance graphs in the plane}, Geombinatorics \textbf{29}(4)
(2020), 137--166; arXiv:2010.12665.

\bibitem{Polymath16}
D.~H.~J. Polymath,
Polymath16: the Hadwiger--Nelson problem, research threads on
D.~G.~Mixon's blog \emph{Short, Fat Matrices} (2018--2021), starting at
\url{https://dustingmixon.wordpress.com/2018/04/14/polymath16-first-thread-simplifying-de-greys-graph/};
project proposal at
\url{https://polymathprojects.org/2018/04/10/polymath-proposal-finding-simpler-unit-distance-graphs-of-chromatic-number-5/}.

\bibitem{ShelahSoifer03}
S.~Shelah and A.~Soifer,
\emph{Axiom of choice and chromatic number of the plane},
J.\ Combin.\ Theory Ser.\ A \textbf{103} (2003), 387--391.

\bibitem{Soifer09}
A.~Soifer,
\emph{The Mathematical Coloring Book},
Springer, New York, 2009.

\bibitem{drattrim}
N.~Wetzler, M.~J.~H. Heule, and W.~A. Hunt,
\emph{DRAT-trim: efficient checking and trimming using expressive
clausal proofs}, in: Theory and Applications of Satisfiability Testing
--- SAT 2014, LNCS 8561, Springer, 2014, 422--429.

\end{thebibliography}

\end{document}
